\chapter{Código fuente de la aplicación móvil}
\label{anexo:app}

Este anexo recoge los fragmentos más representativos del adaptador BLE de la
aplicación móvil React Native. El código completo está disponible en el repositorio
del proyecto (\texttt{mobile-app/src/}).

\section*{C.1 — Constantes NUS y conexión GATT}

\begin{lstlisting}[language=JavaScript, caption={UUIDs del servicio NUS y configuracion de la conexion BLE}]
import { BleManager, Device, Subscription } from 'react-native-ble-plx';

const NUS_SERVICE = '6E400001-B5A3-F393-E0A9-E50E24DCCA9E';
const RX_CHAR     = '6E400002-B5A3-F393-E0A9-E50E24DCCA9E';
const TX_CHAR     = '6E400003-B5A3-F393-E0A9-E50E24DCCA9E';

const MTU                = 128;   // MTU negociado en la conexion
const WRITE_CHUNK_BYTES  = 240;   // troceo de escritura sobre RX
const SCAN_TIMEOUT_MS    = 20_000;
const REQUEST_TIMEOUT_MS = 15_000;
const PING_INTERVAL_MS   = 7_000; // < server timeout (15s) con margen

export class BleAdapter implements IVehicleAdapter {
  private static instance: BleAdapter;
  private readonly manager = new BleManager();
  private device: Device | null = null;
  private txSub: Subscription | null = null;
  private rxBuf = '';
  private queue: Pending[] = [];

  static getInstance(): BleAdapter {
    if (!BleAdapter.instance) BleAdapter.instance = new BleAdapter();
    return BleAdapter.instance;
  }
}
\end{lstlisting}

\begin{lstlisting}[language=JavaScript, caption={Flujo de conexion: GATT, MTU, servicios, suscripcion y autenticacion}]
async connect(deviceId?: string): Promise<void> {
  if (!deviceId) throw new Error('No device selected');

  // 1. Conexion GATT fisica
  const raw = await this.manager.connectToDevice(deviceId,
               { autoConnect: false, timeout: 10000 });

  // 2. Negociacion MTU
  const mtuResult = await raw.requestMTU(MTU);

  // 3. Descubrimiento de servicios y caracteristicas
  await raw.discoverAllServicesAndCharacteristics();

  // 4. Registro de callback de desconexion inesperada
  raw.onDisconnected(() => this.handleUnexpectedDisconnect());
  this.device = raw;

  // 5. Suscripcion a notificaciones TX (RX desde perspectiva del servidor)
  await this.subscribeToTx();

  // 6. Autenticacion por token
  await this.request({ cmd: 'auth', token: '1234' });

  this.startPingInterval();
  this.startActivityMonitor();
}
\end{lstlisting}

\section*{C.2 — Buffer NDJSON y despachador de mensajes}

\begin{lstlisting}[language=JavaScript, caption={Procesamiento del buffer NDJSON recibido por notificaciones BLE}]
private feedBuffer(chunk: string): void {
  this.lastActivityTime = Date.now();
  this.rxBuf += chunk;
  // Separar por lineas completas; la ultima puede estar incompleta
  const lines = this.rxBuf.split('\n');
  this.rxBuf = lines.pop() ?? '';

  for (const line of lines) {
    const t = line.trim();
    if (!t) continue;
    try {
      const parsed = JSON.parse(t) as Record<string, unknown>;
      this.dispatch(parsed);
    } catch {
      // JSON malformado -- descartar fragmento
    }
  }
}

private dispatch(msg: Record<string, unknown>): void {
  // Muestras de monitorizacion continua: enviar a todos los callbacks
  if (msg.type === 'samples') {
    const batch = msg.samples as MonitorSample[];
    batch.forEach((s) => this.sampleCbs.forEach((cb) => cb(s)));
    return;
  }

  // Heartbeat del servidor: responder con ack
  if (msg.type === 'heartbeat') {
    this.sendHeartbeatAck();
    return;
  }

  // Respuesta a comando: resolver o rechazar la promesa en cola
  const pending = this.queue.shift();
  if (!pending) return;
  clearTimeout(pending.timer);

  if (msg.status === 'ok') {
    pending.resolve(msg.data);
  } else {
    pending.reject(new Error((msg.message as string) ?? 'Server error'));
  }
}
\end{lstlisting}

\section*{C.3 — Escritura BLE con chunking y timeout}

\begin{lstlisting}[language=JavaScript, caption={Sistema de peticion-respuesta con timeout y chunking por MTU}]
private request<T = unknown>(cmd: object): Promise<T> {
  return new Promise<T>((resolve, reject) => {
    const json = JSON.stringify(cmd) + '\n';

    // Timeout por si la Raspberry Pi no responde
    const timer = setTimeout(() => {
      const idx = this.queue.findIndex((p) => p.timer === timer);
      if (idx !== -1) this.queue.splice(idx, 1);
      reject(new Error(`Timeout: ${JSON.stringify(cmd)}`));
    }, REQUEST_TIMEOUT_MS);

    this.queue.push({ resolve: resolve as (d: unknown) => void, reject, timer });
    this.writeRx(json).catch((e) => {
      const idx = this.queue.findIndex((p) => p.timer === timer);
      if (idx !== -1) { clearTimeout(timer); this.queue.splice(idx, 1); }
      reject(e);
    });
  });
}

private async writeRx(data: string, withResponse = true): Promise<void> {
  if (!this.device) throw new Error('Not connected');
  const bytes = new TextEncoder().encode(data);

  // Fragmentar en chunks para el caso de payloads grandes
  for (let offset = 0; offset < bytes.length; offset += WRITE_CHUNK_BYTES) {
    const slice = bytes.subarray(offset, offset + WRITE_CHUNK_BYTES);
    let bin = '';
    for (let i = 0; i < slice.length; i++) bin += String.fromCharCode(slice[i]);
    const b64 = btoa(bin);

    if (withResponse) {
      await this.device.writeCharacteristicWithResponseForService(
        NUS_SERVICE, RX_CHAR, b64);
    } else {
      await this.device.writeCharacteristicWithoutResponseForService(
        NUS_SERVICE, RX_CHAR, b64);
    }
  }
}
\end{lstlisting}

\section*{C.4 — Keepalive y monitor de inactividad}

\begin{lstlisting}[language=JavaScript, caption={Ping periodico y desconexion automatica por inactividad}]
private startPingInterval(): void {
  this.pingInterval = setInterval(() => {
    if (!this.device) return;
    // Write sin respuesta: el pong se descarta silenciosamente en dispatch()
    this.writeRx(JSON.stringify({ cmd: 'ping' }) + '\n', false).catch(() => {});
  }, PING_INTERVAL_MS);
}

private startActivityMonitor(): void {
  this.activityMonitorInterval = setInterval(() => {
    if (!this.device) return;
    const elapsed = Date.now() - this.lastActivityTime;
    if (elapsed > CLIENT_TIMEOUT_MS) {
      this.disconnect().catch(() => {});
    }
  }, INACTIVITY_CHECK_MS);
}

private handleUnexpectedDisconnect(): void {
  if (this.pingInterval) clearInterval(this.pingInterval);
  if (this.activityMonitorInterval) clearInterval(this.activityMonitorInterval);
  this.txSub?.remove();
  this.txSub = null;
  this.rxBuf = '';
  this.device = null;
  this.drainQueue(new Error('Device disconnected unexpectedly'));
}
\end{lstlisting}
